\section{Preprocessing and lexicon architecture}
\label{sec:lexicon}

Every zone identified in \S\ref{sec:extraction} passes through a single
\texttt{TextProcessor}: spaCy's \texttt{en\_core\_web\_md} model, with
named-entity recognition and dependency parsing disabled. URLs are stripped,
text is lowercased, tokenised and lemmatised. A token survives only if it is
not a stopword, not punctuation, not whitespace, longer than two characters,
its lemma is not one of 189 corpus-specific stopwords, and it passes a
content-token test: alphanumeric, with at least two alphabetic characters.
That last condition replaced an earlier, purely-alphabetic filter that had
been silently deleting \texttt{co2}, \texttt{co2e} and \texttt{sf6}; it still
excludes bare numerals (\texttt{2023}, \texttt{14001}) and page references
(\texttt{p12}, \texttt{q1}), so digits survive only when attached to
letters. Two representations are carried forward: a lemmatised,
stopword-filtered token string (\texttt{clean\_text}), feeding \SUS, \HEDGE\
and \BREADTH; and a whitespace-collapsed copy of the extracted text with
digits and case intact (\texttt{raw\_text}), the sole input to \QUANT.

\subsection{Vocabulary and corpus in one representation, by construction}

The vocabulary that scores \SUS\ and the lexicon that scores \HEDGE\ are
generated artefacts, not hand-written lists — the generation step is the
methodological point of this section. Each is authored once, in natural
surface form, in a single source file; a build script
(\texttt{scripts/build\_lexicons.py}) passes every entry through the
identical \texttt{TextProcessor} instance applied to the corpus, so that
vocabulary and text are guaranteed to share a representation. The
alternative — a second, hand-lemmatised list kept alongside the raw one —
invites drift: when the corpus pipeline changes, a hand-maintained list
stops tracking it silently. A term whose surface form does not survive
preprocessing becomes a dictionary entry that can never match anything, and
it fails without warning, because a zero count is indistinguishable from the
term's genuine absence.

Two checks in the build step guard against this. A staleness check compares
each source file's modification time against its generated counterpart,
raising an error if the generated file is missing and a warning if the
source has since been edited. A collapse check flags any multi-word term
that reduces to a single word once preprocessed, rather than adding it
silently. The documented case is \emph{paris agreement}, which loses
\emph{paris} to the stopword filter and lemmatises to the bare unigram
\emph{agreement} — a token occurring 15,459 times in the corpus in senses
unrelated to the climate accord (refinancing agreements, employment
terminations, sponsorship deals).

The check is advisory, and the distinction matters in a paper about
undisclosed methodological choices, so we state it exactly rather than
advertise a safeguard the code does not provide. The build script emits a
warning naming each collapsed term and then adds it to the generated
vocabulary regardless; nothing is withheld automatically. Removal is a
separate human decision taken on reading the warning, and it was taken twice:
\emph{paris agreement} and \emph{just transition} were deleted by hand from
\texttt{esg\_terms.txt} and re-entered as literal patterns matched against raw
text at the extraction stage (\S\ref{sec:extraction}), which is why
\emph{agreement} and \emph{transition} are not bare vocabulary entries. A
third collapse was allowed to stand: \emph{due diligence} lemmatises to
\emph{diligence}, which is one of the 154 entries, on the judgement that its
4,330 occurrences in the corpus derive almost entirely from the source phrase.
An automatic rule would have discarded that one too. What the build step
guarantees is that no collapse passes unnoticed; what it does not guarantee is
that a collapsed term stays out of the vocabulary.

\subsection{The ESG vocabulary}

\texttt{esg\_terms.txt} is 154 canonical terms in natural singular form,
organised into seven thematic blocks (Table~\ref{tab:lexicon-blocks}). Run
through the build step, the 154 terms yield 154 distinct lemmatised entries
with no collisions: 85 unigrams, 62 bigrams, 7 trigrams. The resulting
maximum n-gram length, 3, is derived at runtime and supplied as the
vectoriser's \texttt{ngram\_range} upper bound; scikit-learn's
unigrams-only default would otherwise score every multi-word entry — 45\%
of the list — as zero, with no error raised.

\begin{table}[htbp]
\centering
\caption{The 154-term ESG vocabulary, by thematic block.}
\label{tab:lexicon-blocks}
\begin{tabular}{lr}
\toprule
Thematic block & Terms \\
\midrule
Environment: climate and emissions & 14 \\
Environment: energy & 18 \\
Environment: resources and nature & 17 \\
Sustainability and circular economy & 12 \\
Social: people and equity & 34 \\
Governance and ethics & 25 \\
Reporting frameworks and ratings & 34 \\
\midrule
Total & 154 \\
\bottomrule
\end{tabular}
\end{table}

Because \citep{lagasio2024} never states the size of its GRI/SASB-derived
keyword list and publishes only three illustrative terms (\S\ref{sec:original}),
our vocabulary cannot be checked against the original's term for term. This
is where strict replication becomes impossible at the lexicon stage, not
merely inconvenient: any divergence between our \ESGSI\ and the original's
cannot be attributed to a single design choice, since the vocabulary is one
of several elements differing at once and the original's is unknown.

\subsection{The hedge lexicon}

The \HEDGE\ lexicon is built the same way, from a different source.
\texttt{RAW\_LM\_dictionary.csv} supplies the Loughran--McDonald word lists
\citep{loughran2011}, 9,752 rows across seven categories; we keep the four
bearing on hedging and evasive language — Uncertainty, WeakModal,
StrongModal, Constraining, the category names as the source file spells them —
1,245 rows, 1,196 unique lowercase forms.
Loughran--McDonald forms are inflected, and crossing them directly against a
lemmatised corpus recovered only about 46\% of entries. The generated file
is instead the union of each entry's original form and its spaCy lemma,
raising the usable lexicon to 1,263 forms, 67 more than the 1,196 that
entered lemmatisation.

\subsection{External audit of vocabulary coverage}

The vocabulary's coverage was independently audited on 880 zones from four
FY2023 reports (Stellantis, adidas, Danone, Enel). The auditor searched the
zone text with an extended reference lexicon — GRI, SASB, TCFD and ESRS
terminology plus prior literature — by word-boundary regex, excluding any
candidate already in our vocabulary. It identified 147 ESG terms present in
the zone text and absent from ours, 2,970 occurrences in total;
Table~\ref{tab:lexicon-audit} lists the four largest. Because matching was
lexical, with no paraphrase or synonym detection, and the audit covers only
text the extraction stage had already judged ESG-relevant, 147 is a measured
lower bound on the vocabulary's recall, not a full accounting of what it
misses.

\begin{table}[htbp]
\centering
\caption{Largest ESG terms omitted from the vocabulary, external audit.}
\label{tab:lexicon-audit}
\begin{tabular}{llr}
\toprule
Term & Category & Occurrences \\
\midrule
compliance & Governance & 436 \\
social & Social & 275 \\
supply chain & Governance & 241 \\
raw material & Environment & 136 \\
\bottomrule
\end{tabular}
\end{table}

Two readings of the shortfall point in opposite directions, and the one that
flatters us is wrong. By term count it looks sector-specific: 64 of the 147
appear in only one of the four audited firms. Weighted by occurrence the
impression reverses. Those 64 single-firm terms account for 450 of the 2,970
occurrences (15\%), while the 22 terms present in \emph{all four} firms —
including every row of Table~\ref{tab:lexicon-audit} — account for 1,573
(53\%). The gap is dominated by general ESG vocabulary missing everywhere the
vocabulary was built to look, not by sector idiosyncrasy, so single-firm
concentration is not a basis for discounting it.

The sector concern was in any case an artefact of the sample that measured it.
The 147 collapse to 143 distinct candidate strings once near-duplicate surface
forms are merged, and every one of them has since been re-measured over all 343
reports — occurrences, documents, companies, countries, and the share held by
the largest single company. On that base 103 are broadly dispersed — 102 of
which reached the vocabulary, \emph{power purchase agreement} having been lost
to a generation-script bug after the audit rejected the bare \emph{ppa} in its
favour — 18 are frequent but polysemous, and 22 are concentrated enough in one
industry to be called sectoral, carrying 2.0\% of the candidates' occurrence
mass. Four companies were too few to tell a sectoral term from a common one.
The dispersed and the polysemous classes were incorporated into an expanded
289-entry vocabulary; the 22 sectoral terms were not, a
choice \S\ref{sec:limitations} records together with its measured effect. The
specification reported throughout this paper remains the 154 terms of
Table~\ref{tab:lexicon-blocks}, and the expansion is a robustness check on it
(\S\ref{sec:robustness-lexicon}).

Two candidates carry roughly a third of the occurrence mass between them.
Counted as whole words over the full text of all 343 reports rather than inside
the four audited ones, \emph{compliance} occurs 55,780 times and \emph{social}
48,198. Neither is admissible as a bare term, and they fail differently
(Table~\ref{tab:lexicon-context}). Of \emph{social}'s
occurrences 60.3\% sit in unambiguously ESG context and only 8.7\%
unambiguously outside it, concentrated in three enumerable strings —
\emph{social security}, \emph{social media}, \emph{social network} — so on the
evidence the right treatment of \emph{social} is an exception list. The
vocabulary architecture supplies none: a bare entry in a TF-IDF vocabulary
cannot carry a negative lookahead, and no exclusion mechanism exists at any
stage. What shipped for \emph{social} is therefore the opposite of what its own
analysis recommends — an inclusion list of eight collocations, admitting only
the uses they name and discarding the long tail of valid ones the analysis
identified as the majority of its mass. \emph{Compliance} inverts the evidence:
41.8\% are unambiguously ESG, 8.7\% are accounting or audit usage sitting
almost entirely in the auditor's report and the notes to the financial
statements, and the largest block, 44.5\%, carries no marker either way, so for
that term an inclusion list of fifteen collocations is what the evidence
supports as well as what was built. \emph{Proxy} was adjudicated on the same
evidence, against \emph{proxy hedging}.

These shares are measured over the full text of each report, not over the
extracted zones, and the two bases differ in exactly the direction that
matters: extraction has already discarded most of the notes and the auditor's
report, which is where the accounting sense of \emph{compliance} lives. The
false-positive rate inside the pipeline is therefore lower than
Table~\ref{tab:lexicon-context} shows, which bounds the problem from above and
does not estimate it.

What the dispersion analysis retires is the sector-bias objection, which is now
measured, and not the recall bound, which is unchanged: the adopted vocabulary
incorporates none of the 147, and every figure in this paper is computed on the
154 entries. The bound stands exactly where \S\ref{sec:limitations} leaves it.

\begin{table}[htbp]
\centering
\caption{Context of the two largest audit candidates, classified by surrounding
vocabulary over the full text of all 343 reports. Shares of all occurrences of
each term, per cent.}
\label{tab:lexicon-context}
\begin{tabular}{lrr}
\toprule
Context & \emph{social} & \emph{compliance} \\
\midrule
Unambiguously ESG & 60.3 & 41.8 \\
Mixed & 10.6 & 3.8 \\
No marker either way & 20.4 & 44.5 \\
Not ESG: payroll, tax, communications & 7.1 & --- \\
Not ESG: accounting or audit & 1.6 & 8.7 \\
Not ESG: other & --- & 1.2 \\
\midrule
Total & 100.0 & 100.0 \\
\bottomrule
\end{tabular}
\end{table}
